\documentclass{article}
\usepackage{amsmath}
\usepackage{amssymb}
\usepackage{amsthm}
\usepackage{cite}

\title{Formal Definition of Multi-Agent Coordination Experimental Framework}
\author{}
\date{}

\begin{document}
\maketitle

\section{System Definition}

\subsection{Agent Set}
Let $\mathcal{A} = \{a_1, a_2, a_3, a_u\}$ be a finite set of four agents where:
\begin{itemize}
    \item $a_1, a_2, a_3$ are \textbf{defined agents} with observable behavioral patterns
    \item $a_u$ is an \textbf{undefined agent} with unspecified behavioral parameters
\end{itemize}

\subsection{State Space}
The system operates within state space $\Omega$ where each state $\omega \in \Omega$ represents:
$$\omega = (x_1(t), x_2(t), x_3(t), x_u(t), e(t))$$
where $x_i(t) \in \mathbb{R}^3$ represents spatial coordinates of agent $i$ at time $t$, and $e(t)$ represents environmental conditions.

\subsection{Action Space}
Each agent $a_i$ has access to action set $\mathcal{S}_i$ where:
$$\mathcal{S}_i = \{\text{stationary}, \text{locomotion}_{\vec{v}}, \text{signal}_s\}$$
with $\vec{v} \in \mathbb{R}^3$ representing velocity vectors and $s$ representing communication signals.

\section{Information Structure}

\subsection{Observable Variables}
Each agent $a_i$ has access to information set $\mathcal{I}_i(t)$ consisting of:
\begin{align}
\mathcal{I}_i(t) = \{&\text{spatial positions of other agents within sensory range} \nonumber \\
&\text{velocity vectors of observable agents} \nonumber \\
&\text{environmental stimuli within detection threshold}\}
\end{align}

\subsection{Communication Protocol}
Agents may exchange information through signal set $\mathcal{M}$ governed by:
$$\mathcal{M}: \mathcal{A} \times \mathcal{A} \rightarrow \{0, 1\}^k$$
where $k$ represents signal dimensionality and transmission success is probabilistic.

\section{Behavioral Dynamics}

\subsection{Response Functions}
Agent behavior is governed by response function $f_i: \mathcal{I}_i(t) \rightarrow \mathcal{S}_i$ where:
$$a_i(t+1) = f_i(\mathcal{I}_i(t), \theta_i)$$
with $\theta_i$ representing agent-specific parameters (potentially unknown for $a_u$).

\subsection{Coordination Mechanism}
The system exhibits emergent coordination through:
\begin{theorem}[Behavioral Convergence]
Given initial divergent states, agents achieve synchronized behavioral patterns without explicit coordination protocols.
\end{theorem}

\section{Environmental Parameters}

\subsection{Spatial Constraints}
The experimental space $\mathcal{E} \subset \mathbb{R}^3$ may contain:
\begin{itemize}
    \item Physical barriers $B \subset \mathcal{E}$
    \item Visibility constraints $V(\mathcal{E}) \rightarrow [0,1]$
    \item Accessibility regions $\mathcal{A}cc(\mathcal{E}) \subseteq \mathcal{E}$
\end{itemize}

\subsection{External Entities}
Additional entities $\mathcal{X} = \{x_1, x_2, ..., x_n\}$ may be introduced with:
$$\mathcal{X}_i: \{\text{autonomous}, \text{constrained}, \text{static}\}$$
where classification affects interaction potential with agent set $\mathcal{A}$.

\section{Measurement Framework}

\subsection{Observables}
The experimental framework measures:
\begin{align}
\text{Coordination Index} &= \frac{\sum_{i=1}^{3} \|\vec{v}_i(t) - \bar{\vec{v}}(t)\|^2}{\sum_{i=1}^{3} \|\vec{v}_i(t)\|^2} \\
\text{Response Latency} &= t_{\text{response}} - t_{\text{stimulus}} \\
\text{Behavioral Entropy} &= -\sum_{s \in \mathcal{S}} P(s) \log P(s)
\end{align}

\subsection{Invariance Properties}
\begin{theorem}[Classification Invariance]
Coordination measures remain invariant under relabeling of entities within environmental space.
\end{theorem}

\section{Experimental Protocols}

\subsection{Control Conditions}
Baseline measurements require:
\begin{itemize}
    \item Isolation of agent behavioral parameters
    \item Elimination of external stimuli
    \item Standardized initial conditions
\end{itemize}

\subsection{Perturbation Analysis}
Systematic introduction of variables $\Delta_k$ where:
$$\Delta_k \in \{\text{entity introduction}, \text{environmental modification}, \text{information constraint}\}$$
enables measurement of system robustness to external modifications.

\section{Theoretical Implications}

\subsection{Game-Theoretic Framework}
The system operates without explicit payoff functions, suggesting:
$$\nexists u_i: \Omega \rightarrow \mathbb{R} \text{ such that } \max_{s_i} \mathbb{E}[u_i(\omega(s_i, s_{-i}))]$$
defines agent behavior, contradicting standard game-theoretic assumptions \cite{nash1950equilibrium}.

\subsection{Information-Theoretic Analysis}
Coordination emerges despite information asymmetries, violating common knowledge requirements for rational coordination \cite{aumann1976agreeing}.

\bibliographystyle{plain}
\begin{thebibliography}{2}
\bibitem{nash1950equilibrium} Nash, J. (1950). Equilibrium points in n-person games. \textit{Proceedings of the National Academy of Sciences}, 36(1), 48-49.
\bibitem{aumann1976agreeing} Aumann, R. J. (1976). Agreeing to disagree. \textit{The Annals of Statistics}, 4(6), 1236-1239.
\end{thebibliography}

\end{document}
